Absolute aggregation raised PSNR by 0.487 dB at the low threshold and by 1.089 dB at the high threshold, an interaction of about 0.602 dB. It also increased count at each fixed threshold. The absolute low-threshold arm reached 31.084 dB and 699,645 Gaussians, with competition for slots at all 54 events. The selected high-threshold arm gave up 0.327 dB and saved 247,236 primitives (35.3\%). It still improved on B00's 30.125 dB by 0.632 dB with 45.8\% fewer primitives. The high-threshold signed arm grew to only 143,086 Gaussians and fell below B00 on all three RGB metrics. These contrasts measure the gradient eligibility policy under fixed split and ranking settings. The low threshold also activates the cap, so the threshold contrast includes that downstream effect.

The earlier A04 factorial supports a contribution from Long-Axis Split and replacing clones with splits under signed gradients (Table~\ref{tab:split-ablation}). An equivalent factorial under A05-Abs would measure whether the split effects change with absolute gradients [\hyperlink{evE3}{E3}].

At 700,000, every eligible parent split in all three absolute/high-threshold ranking arms. At 400,000, competition occurred in 32 or 33 events. Removing the RGB-error cue changed PSNR by $-0.018$ dB and removing the edge cue changed it by $+0.033$ dB, at nearly equal counts. Both removals slightly increased LPIPS. The full configuration remains the fixed reference because it is the design carried forward from development and has the lowest LPIPS in this binding-cap trio. The edge-free variant has the highest PSNR here. A single Room seed and opposing metric directions leave the general value of the cues unresolved [\hyperlink{evE8}{E8}].
